%!TEX root=../../note
\subsection{The Algebraic and Order Properties of $\R$}
The real numbers $\R$ form a \emph{complete ordered field}. We build up this
description in three stages: the field (algebraic) axioms, the order axioms, and
the completeness axiom.

\begin{definition}[Field]
    A \textbf{field} is a set $F$ equipped with two binary operations
    \[
        + : F \times F \to F,
        \qquad
        \cdot : F \times F \to F,
    \]
    satisfying the following axioms for all $a,b,c \in F$.

    \begin{enumerate}[label=\textbf{F\arabic*.}, leftmargin=2.5em]
        \item \textbf{Associativity of addition.}
        \[
            a + (b + c) = (a + b) + c.
        \]

        \item \textbf{Commutativity of addition.}
        \[
            a + b = b + a.
        \]

        \item \textbf{Additive identity.}
        There exists an element $0 \in F$ such that
        \[
            a + 0 = a.
        \]

        \item \textbf{Additive inverses.}
        For every $a \in F$, there exists an element $-a \in F$ such that
        \[
            a + (-a) = 0.
        \]

        \item \textbf{Associativity of multiplication.}
        \[
            a \cdot (b \cdot c) = (a \cdot b) \cdot c.
        \]

        \item \textbf{Commutativity of multiplication.}
        \[
            a \cdot b = b \cdot a.
        \]

        \item \textbf{Multiplicative identity.}
        There exists an element $1 \in F$, with $1 \neq 0$, such that
        \[
            a \cdot 1 = a.
        \]

        \item \textbf{Multiplicative inverses.}
        For every $a \in F$ with $a \neq 0$, there exists an element $a^{-1} \in F$ such that
        \[
            a \cdot a^{-1} = 1.
        \]

        \item \textbf{Distributive law.}
        \[
            a \cdot (b + c) = a \cdot b + a \cdot c.
        \]
    \end{enumerate}
\end{definition}

\begin{remark}
    The rational numbers $\Q$ also form a field.
\end{remark}

\begin{theorem}
    \begin{enumerate}[label = \alph*)]
        \item If $z, a \in \R$ with $z + a = a$, then $z = 0$.
        \item If $u \in \R$ and $b \neq 0$ with $u \cdot b = b$, then $u = 1$.
        \item If $a \in \R$, then $a \cdot 0 = 0$.
    \end{enumerate}
\end{theorem}

\begin{proof}
    \begin{enumerate}[label = \alph*)]
        \item
        \[
            z = z + 0 = z + (a - a) = (z + a) + (-a) = a + (-a) = 0.
        \]
        \item Multiplying $u \cdot b = b$ by $b^{-1}$,
        \[
            u = u \cdot 1 = u \cdot \paren{b \cdot b^{-1}} = (u \cdot b) \cdot b^{-1} = b \cdot b^{-1} = 1.
        \]
        \item
        \[
            a\cdot0=a\cdot(0+0)=a\cdot0+a\cdot0.
        \]
        Adding $-(a\cdot0)$ to both sides, and using associativity, gives
        $0=a\cdot0$.
    \end{enumerate}
\end{proof}

\begin{definition}[The Order Properties of $\R$]
    There is a nonempty subset $P$ of $\R$, called the set of \textbf{positive real
    numbers}, that satisfies the following properties.
    \begin{enumerate}[label=\textbf{O\arabic*.}, leftmargin=3em]
        \item \textbf{Trichotomy law.}
        For every $a \in \R$, exactly one of the following holds:
        \begin{equation*}
            a = 0, \qquad a \in P, \qquad -a \in P.
        \end{equation*}

        \item \textbf{Closure under addition.}
        If $a, b \in P$, then $a + b \in P$.

        \item \textbf{Closure under multiplication.}
        If $a, b \in P$, then $a \cdot b \in P$.
    \end{enumerate}
\end{definition}

\begin{definition}[Order relation]
    Let $a, b \in \R$.
    \begin{enumerate}
        \item If $a - b \in P$, then we write $a > b$ or $b < a$.
        \item If $a - b \in P \cup \set{0}$, then we write $a \geq b$ or $b \leq a$.
    \end{enumerate}
\end{definition}

\begin{remark}
    $\Q$ is also ordered, so we are still missing something: completeness.
\end{remark}

\begin{definition}[Completeness Axiom]
    If $A \subseteq \R$ is nonempty and bounded above, then $A$ has a least upper
    bound in $\R$.
\end{definition}

\begin{example}
    Prove that if $a \in \R$ and $a \neq 0$, then $a^2 > 0$.
\end{example}

\begin{proof}
    Since $a \neq 0$, the trichotomy law gives two cases.

    \textbf{Case 1: $a \in P$.} By closure under multiplication,
    \begin{equation*}
        a^2 = a \cdot a \in P, \qquad\text{so } a^2 > 0.
    \end{equation*}

    \textbf{Case 2: $a \notin P$.} Since $a \neq 0$, trichotomy gives $-a \in P$, so by closure under multiplication $(-a)(-a) \in P$. It remains to show that $a^2 = (-a)(-a)$:
    \begin{align*}
        a + (-a) = 0
            &\implies (-a)\paren{a + (-a)} = (-a) \cdot 0 = 0 \\
            &\implies (-a)(a) + (-a)(-a) = 0 \\
            &\implies (-a)(-a) = -\paren{(-a)a} = -(-a^2) = a^2.
    \end{align*}
    Hence $a^2 = (-a)(-a) \in P$, so $a^2 > 0$. In either case $a^2 > 0$.
\end{proof}

\subsubsection{The Absolute Value}
\begin{definition}
    Given $a \in \R$, we define the \textbf{absolute value} $\abs{a}$ by
    \[
        \abs{a} =
        \begin{cases}
            a, & a \geq 0, \\
            -a, & a < 0.
        \end{cases}
    \]
\end{definition}
\begin{theorem}
    \begin{enumerate} [label = (\alph*)]
        \item $\abs{ab} = \abs{a}\abs{b}$ for all $a, b \in \R$.
        \item $\abs{a}^2 = a^2$ for all $a \in \R$.
        \item For $a \in \R$ and $c>0$,
        \[\abs{a} \leq c \iff -c \leq a \leq c.\]
        \item For every $a \in \R$,
        \[-\abs{a} \leq a \leq \abs{a}.\]

    \end{enumerate}

\end{theorem}

\begin{proof}
    The definition immediately gives $\abs{-a}=\abs{a}$: if $a\geq0$,
    then $-a\leq0$, and if $a<0$, then $-a>0$. It also gives
    \begin{equation*}
        -\abs{a}\leq a\leq\abs{a},
    \end{equation*}
    proving (d). In the two cases $a\geq0$ and $a<0$, respectively,
    $\abs{a}^2=a^2$ is immediate; this proves (b).

    For (a), there are four sign combinations. If $a,b\geq0$, then
    $ab\geq0$, so $\abs{ab}=ab=\abs a\abs b$. If $a\geq0>b$, then
    $ab\leq0$, so $\abs{ab}=-ab=a(-b)=\abs a\abs b$. The case
    $b\geq0>a$ is the same with the roles reversed. Finally, if $a,b<0$,
    then $ab>0$, and
    \begin{equation*}
        \abs{ab}=ab=(-a)(-b)=\abs a\abs b.
    \end{equation*}

    Finally let $c>0$. Suppose first that $\abs a\leq c$. If $a\geq0$,
    then $\abs a=a$, so $a\leq c$; also $-c<0\leq a$. If $a<0$, then
    $\abs a=-a$, so $-a\leq c$, equivalently $-c\leq a$; also $a<0<c$.
    Hence $-c\leq a\leq c$ in either case.

    Conversely, suppose $-c\leq a\leq c$. If $a\geq0$, then
    $\abs a=a\leq c$. If $a<0$, multiplying $-c\leq a$ by $-1$ reverses
    the inequality and gives $-a\leq c$, hence $\abs a\leq c$. This proves (c).
\end{proof}

\begin{proposition}
    [Triangle Inequality]
    If $a, b \in \R$, then $\abs{a + b } \leq \abs{a } + \abs{b }$.
\end{proposition}

\begin{proof}
    Let $a, b \in \R$. We know that
    \begin{align*}
        -\abs{a } \leq a \leq \abs{a}. \\
        -\abs{b } \leq b \leq \abs{b}.
    \end{align*}

    Thus,
    \begin{align*}
        &-\abs{a } - \abs{b } \leq a + b \leq \abs{a } + \abs{ b} \\
        &\iff - \paren{\abs{a } + \abs{b}} \leq a + b \leq \abs{a} + \abs{b}
    \end{align*}

    If $\abs{a}+\abs{b}=0$, then the displayed bounds give $a+b=0$, so the
    claim is immediate. Otherwise, part (c) of the preceding theorem gives
    \begin{equation*}
        \abs{a + b } \leq \abs{a} + \abs{b}.
    \end{equation*}
\end{proof}

\begin{corollary}
    \begin{enumerate}
        \item $\abs{\abs{a } - \abs{b}} \leq \abs{a - b}. $
        \item $\abs{a - b } \leq \abs{a } + \abs{b}$
    \end{enumerate}
\end{corollary}

\begin{proof}[Proof of (1)]
    Let $a, b \in \R$. We know
    \begin{align*}
        a = a - b + b &\implies \abs{a} = \abs{(a - b) + b} \leq \abs{a - b} + \abs{b} \\
        &\implies \abs{a } - \abs{b } \leq \abs{a - b}
    \end{align*}

    \begin{align*}
        b = b - a + a &\implies \abs{b } = \abs{(b - a) + a } \leq \abs{b - a } + \abs{a }\\
        &\implies \abs{b } - \abs{a } \leq \abs{b-a}
            = \abs{-(a-b)}=\abs{a-b}. \\
        &\implies \abs{a } - \abs{b } \geq - \abs{a - b}
    \end{align*}

    Thus,
    \begin{equation*}
        -\abs{a - b } \leq \abs{a } - \abs{b } \leq \abs{a - b }
    \end{equation*}

    Therefore,
    \begin{equation*}
        \abs{\abs{a } - \abs{b }} \leq \abs{a - b}
    \end{equation*}
\end{proof}

\begin{proof}[Proof of (2)]
    Let $a, b \in \R$. We know
    \begin{align*}
        \abs{a - b} = \abs{a + \paren{-b }} \leq \abs{a} + \abs{-b }
        = \abs{a } + \abs{b}.
    \end{align*}
\end{proof}
